\documentclass{article}

% NeurIPS 2026 Submission Template
% Pass nonatbib if natbib package clashes
\usepackage[preprint]{neurips_2026}

\usepackage[utf8]{inputenc} % allow utf-8 input
\usepackage[T1]{fontenc}    % use 8-bit T1 fonts
\usepackage{hyperref}       % hyperlinks
\usepackage{url}            % simple URL typesetting
\usepackage{booktabs}       % professional-quality tables
\usepackage{amsfonts}       % blackboard math symbols
\usepackage{nicefrac}       % compact symbols for 1/2, etc.
\usepackage{microtype}      % microtypography
\usepackage{xcolor}         % colors
\usepackage{graphicx}       % images
% Fallback macro: embeds real image if file exists, or renders placeholder box if pending
\newcommand{\safeincludegraphics}[2][]{%
  \IfFileExists{#2}{\includegraphics[#1]{#2}}{\fbox{\parbox[c][3.2cm][c]{\linewidth}{\centering \small \texttt{[Figure pending generation: #2]}}}}%
}
\usepackage{amsmath}        % math symbols
\usepackage{subcaption}     % subfigures
\usepackage{float}          % figure positioning

\title{EE-5102 / CS-6304: Advanced Topics in Machine Learning\\Programming Assignment 1: Learning Beyond IID and Closed-Set Assumptions}

\author{%
Abdullah Fayyaz \\
Department of Computer Science / Electrical Engineering \\
Lahore University of Management Sciences \\
\texttt{28100069@lums.edu.pk} \\
\vspace{1mm} \\
}

\begin{document}

\maketitle

\begin{abstract}
Standard supervised learning generally assumes that training and test data are drawn from the same distribution and share the same label space. These assumptions are frequently violated in deployment. This paper investigates representation learning beyond these assumptions by evaluating architectural inductive biases, unsupervised domain adaptation and generalization under appearance shifts, and open-set recognition. Our experiments reproduced the finding that Vision Transformers exhibit higher shape bias than CNNs. Under appearance shifts, our experiments found no direct link between domain invariance and target recognition; the data suggested a negative relationship between domain alignment pressure and accuracy. Finally, we demonstrated that high closed-set accuracy did not always prevent false acceptance of unknown classes. Code and configuration files to reproduce all experiments are available at: \url{https://github.com/FayyazAbdullah2005/PA1-ATML}.
\end{abstract}

%%%%%%%%%%%%%%%%%%%%%%%%%%%%%%%%%%%%%%%%%%%%%%%%%%%%%%%%%%%%%%%%%%%%%%%%%%%%%%%
\section{Introduction and Research Questions}
\label{sec:introduction}

Across the four tasks, this paper investigates the following core research questions:
\begin{enumerate}
    \item \textbf{Representations (Task 1):} How do convolutional networks and Vision Transformers differ in their reliance on shape, texture, color, and spatial layout?
    \item \textbf{Domain Adaptation (Task 2):} Does aligning feature distributions with an unlabeled target domain improve target recognition, or can it induce negative transfer?
    \item \textbf{Domain Generalization (Task 3):} Without access to target data, does domain invariance or local parameter stability improve generalization to an unseen domain?
    \item \textbf{Open-Set Recognition (Task 4):} How can a classifier identify inputs from unknown classes rather than confidently assigning every input to one of its known classes?
\end{enumerate}

%%%%%%%%%%%%%%%%%%%%%%%%%%%%%%%%%%%%%%%%%%%%%%%%%%%%%%%%%%%%%%%%%%%%%%%%%%%%%%%
\section{Task 1: Inductive Biases and Feature Representations}
\label{sec:task1}

\subsection{Methodology}
Since the three models have shown high accuracies on large datasets (e.g., ImageNet), we expect to observe high scores on other similar datasets. We measure performance using Top-1 Accuracy, Macro-F1, and Mean Max Confidence. To evaluate the three different models (ResNet-50, ViT-B/16, and OpenCLIP), the STL-10 dataset was used with an 80/20 training/validation split and a seed of 6304. For the three models, the following backbones were frozen and used:
\begin{enumerate}
    \item torchvision ResNet-50 with \texttt{ResNet50\_Weights.IMAGENET1K\_V2}
    \item torchvision ViT-B/16 with \texttt{ViT\_B\_16\_Weights.IMAGENET1K\_V1}
    \item OpenCLIP ViT-B-32 with \texttt{pretrained='openai'}
\end{enumerate}
The linear classifier head for each backbone was then trained using the AdamW optimizer (learning rate $10^{-3}$, weight decay $10^{-4}$, and early stopping patience of 5). Along with these three models, OpenCLIP Zero-Shot classification was also recorded using the prompt \texttt{"a photo of a \{class\}"}.

Since experiments conducted by \citet{geirhos2019imagenet} on texture bias in CNNs showed that grayscaling images did not result in a significant drop in performance for color interventions, we similarly expect little or no difference in performance for ResNet here. Since ViT and OpenCLIP learn global features and are trained on large datasets, we expect these models to also show only small drops in performance. We quantify performance changes using accuracy change ($\Delta \text{Acc}$) and prediction consistency. To test color bias, images were passed through a grayscale filter as the first test and a $180^\circ$ Hue rotation as a second test (evaluated across the 500 clean test images, seed 6304).

\citet{geirhos2019imagenet} found a strong bias towards local textures in CNNs and \citet{raghu2021vision} found that ViTs produce more uniform representations across layers, so we expect CNNs to be more texture-biased and ViTs to be more shape-biased. Evaluated metrics are Shape Bias (\%) and Coverage (\%). A total of 220 cue-conflict images were created using AdaIN style transfer ($\alpha=0.8$) across 5 bidirectional class pairs, filtered by pre-evaluation visual rejection criteria ($\text{SSIM} \ge 0.35$, Sobel edge correlation $\ge 0.40$). Model predictions were categorized into shape, texture, or third-class (other) decisions.

Pixel-level translations of $\delta \in \{0, 8, 16, 32\}$ were applied in four cardinal directions (North, South, East, West) using reflection padding to the 500 test images. The models were evaluated on these images in order to test whether predictions remain stable. Since all three models were trained with random crop augmentations during pretraining, we hypothesize that the models will remain fairly stable under small shifts.

To test the significance of patch structure for each model, each image was divided into a $4 \times 4$ pixel-space grid (yielding 16 square tiles of $56 \times 56$ pixels) and scrambled using a deterministic non-identity patch permutation with seed 6304 across all 500 test images. Top-1 Accuracy was computed along with accuracy change ($\Delta \text{Acc}$) and Consistency scores. Our working hypothesis is that models biased towards local features will perform better in this task (acting as a ``bag of features''), since this intervention destroys global layout while preserving local patch evidence.

In order to examine whether the interventions altered the internal representations even when predictions remained stable, we paired each transformed image with its clean counterpart and measured the cosine stability ($I_T$) of each backbone's representation. We also fit 2D UMAP projections on combined clean and transformed features using seed 6304.

\subsection{Results}

\paragraph{Clean Baseline Evaluation}
\begin{table}[H]
    \centering
    \caption{\textbf{Clean Baseline Evaluation}: Comparison of linear classifier heads (frozen backbones) and zero-shot OpenCLIP on the balanced 500-image test subset (seed 6304).}
    \label{tab:task1_clean_baseline}
    \begin{tabular}{lccc}
        \toprule
        \textbf{Model / Backbone} & \textbf{Top-1 Accuracy (\%)} & \textbf{Macro-F1} & \textbf{Mean Max Confidence} \\
        \midrule
        ResNet-50 (Linear Head) & 97.40 & 0.9740 & 0.9065 \\
        ViT-B/16 (Linear Head) & 98.00 & 0.9800 & 0.9713 \\
        OpenCLIP ViT-B-32 (Linear Head) & 97.40 & 0.9739 & 0.1397 \\
        OpenCLIP ViT-B-32 (Zero-Shot) & 68.20 & 0.6509 & 0.8977 \\
        \bottomrule
    \end{tabular}
\end{table}
The results show high accuracies for all three models ($>95\%$) with ViT having the highest accuracy, confidence, and F1-score. ResNet and OpenCLIP (Linear Head) showed very similar accuracy and F1 scores, though ResNet showed higher confidence. OpenCLIP Zero-Shot classification showed significantly lower accuracy at $68.20\%$ and much higher confidence ($0.8977$) compared to OpenCLIP's confidence with a linear probe ($0.1397$). These results are expected, since zero-shot classification is performed using a generic prompt without any fine-tuning.

\paragraph{Chromatic Reliance vs. Shape and Texture Bias}
\begin{table}[H]
    \centering
    \caption{\textbf{Impact of Color Interventions}: Top-1 accuracy change ($\Delta$) and prediction consistency relative to clean images.}
    \label{tab:task1_color}
    \begin{tabular}{lcccc}
        \toprule
        & \multicolumn{2}{c}{\textbf{Grayscale}} & \multicolumn{2}{c}{\textbf{180$^\circ$ Hue Rotation}} \\
        \cmidrule(lr){2-3} \cmidrule(lr){4-5}
        \textbf{Model} & \textbf{$\Delta$ Acc (\%)} & \textbf{Consistency (\%)} & \textbf{$\Delta$ Acc (\%)} & \textbf{Consistency (\%)} \\
        \midrule
        ResNet-50 & -2.40 & 96.80 & -1.20 & 97.00 \\
        ViT-B/16 & -1.60 & 96.80 & -1.80 & 96.20 \\
        OpenCLIP ViT-B-32 (Probe) & -2.80 & 95.40 & -3.40 & 94.80 \\
        OpenCLIP ViT-B-32 (Zero-Shot) & -1.40 & 94.60 & -3.20 & 90.40 \\
        \bottomrule
    \end{tabular}
\end{table}
The results show a mild drop in performance for each model in accuracy and consistency scores. The results are consistent with our original hypothesis with the exception of OpenCLIP. The results suggest that OpenCLIP is the most susceptible to Hue rotation as it showed the largest drops in accuracy ($-3.40\%$ probe, $-3.20\%$ zero-shot) and consistency scores ($90.40\%$ zero-shot). An interpretation of these results is that rotating hue alters the color patterns that OpenCLIP's text and image embeddings rely on, which causes the top-1 prediction to change.

\begin{figure}[H]
    \centering
    \begin{subfigure}[b]{0.48\linewidth}
        \centering
        \safeincludegraphics[width=\linewidth]{figures/task1/cue_conflict_samples.png}
        \caption{Stylized Cue-Conflict Examples}
        \label{fig:cue_conflict_samples}
    \end{subfigure}
    \hfill
    \begin{subfigure}[b]{0.48\linewidth}
        \centering
        \safeincludegraphics[width=\linewidth]{figures/task1/rejection_rule_examples.png}
        \caption{Visual Inspection \& Rejection Examples}
        \label{fig:rejection_rule_examples}
    \end{subfigure}
    \caption{\textbf{AdaIN Cue-Conflict Generation}: (a) Bidirectional shape/texture conflict examples across chosen class pairs; (b) Representative accepted vs. rejected stylizations based on defined visual criteria.}
    \label{fig:task1_stylization}
\end{figure}

\begin{table}[H]
    \centering
    \caption{\textbf{Shape vs. Texture Decisions and Bias Metrics}: Decision counts, Shape Bias (\%), and Coverage (\%) across models on 220 valid cue-conflict images.}
    \label{tab:task1_shape_bias}
    \begin{tabular}{lccccc}
        \toprule
        \textbf{Model} & \textbf{$N_{\text{shape}}$} & \textbf{$N_{\text{texture}}$} & \textbf{$N_{\text{other}}$} & \textbf{Shape Bias (\%)} & \textbf{Coverage (\%)} \\
        \midrule
        ResNet-50 & 119 & 32 & 69 & 78.81 & 68.64 \\
        ViT-B/16 & 163 & 10 & 47 & 94.22 & 78.64 \\
        OpenCLIP ViT-B-32 (Probe) & 148 & 15 & 57 & 90.80 & 74.09 \\
        OpenCLIP ViT-B-32 (Zero-Shot) & 103 & 32 & 85 & 76.30 & 61.36 \\
        \bottomrule
    \end{tabular}
\end{table}

\begin{figure}[H]
    \centering
    \safeincludegraphics[width=0.9\linewidth]{figures/task1/cue_conflict_failures.png}
    \caption{\textbf{Qualitative Cue-Conflict Predictions}: Representative examples highlighting model agreements, disagreements (shape preference vs. texture preference), and third-class failure modes.}
    \label{fig:task1_conflict_cases}
\end{figure}
For cue conflicts, the results show that ResNet-50 exhibits relatively lower shape bias ($78.81\%$) compared to OpenCLIP Probe ($90.80\%$) and ViT ($94.22\%$), choosing texture 32 times. ViT shows the highest coverage out of all models at $78.64\%$ followed by OpenCLIP Probe ($74.09\%$) and then ResNet-50 ($68.64\%$), with OpenCLIP Zero-Shot displaying the lowest coverage ($61.36\%$) and lowest shape bias ($76.30\%$). These results are consistent with our original hypothesis and show that OpenCLIP Zero-Shot displays the highest confusion ($N_{\text{other}}=85$) when predicting labels for cue-conflict images. The reason for these results is that CNNs do in fact have a higher texture bias due to local convolutions, while the self-attention mechanism allows ViT to learn global boundaries better. Furthermore, OpenCLIP Probe also exhibits high shape bias due to its Vision Transformer backbone; the lower shape bias and coverage in Zero-Shot can be explained by the models' high confusion when textures conflict with text prompt semantics, as illustrated in Figure~\ref{fig:task1_conflict_cases}.

\paragraph{Spatial Invariance: Translation vs. Patch Permutation}
\begin{figure}[H]
    \centering
    \begin{subfigure}[b]{0.48\linewidth}
        \centering
        \safeincludegraphics[width=\linewidth]{figures/task1/translation_accuracy.png}
        \caption{Top-1 Accuracy vs. Displacement}
        \label{fig:translation_acc}
    \end{subfigure}
    \hfill
    \begin{subfigure}[b]{0.48\linewidth}
        \centering
        \safeincludegraphics[width=\linewidth]{figures/task1/translation_consistency.png}
        \caption{Prediction Consistency vs. Displacement}
        \label{fig:translation_consistency}
    \end{subfigure}
    \caption{\textbf{Translation Sensitivity Analysis}: (a) Top-1 accuracy and (b) prediction consistency $\text{Consistency}(\delta)$ plotted against spatial pixel displacements $\delta \in \{0, 8, 16, 32\}$ averaged across four cardinal directions.}
    \label{fig:task1_translation_curves}
\end{figure}
The results show near-zero changes in accuracy and prediction consistency for all models under translation ($>96.5\%$ consistency even at $\delta=32$). This suggests that spatial displacement has a negligible effect on model accuracy, though OpenCLIP ViT-B/32 showed a slightly larger drop ($97.05\%$ probe consistency) because a 32-pixel shift spans an entire patch stride on its $32 \times 32$ grid.

\begin{table}[H]
    \centering
    \caption{\textbf{Patch Structure Disruption ($4 \times 4$ Shuffling)}: Performance drop and prediction consistency relative to clean images.}
    \label{tab:task1_patch_shuffle}
    \begin{tabular}{lccc}
        \toprule
        \textbf{Model} & \textbf{Shuffled Top-1 Acc (\%)} & \textbf{$\Delta$ Acc (\%)} & \textbf{Consistency (\%)} \\
        \midrule
        ResNet-50 & 91.40 & -6.00 & 92.00 \\
        ViT-B/16 & 90.80 & -7.20 & 91.00 \\
        OpenCLIP ViT-B-32 (Probe) & 80.60 & -16.80 & 82.00 \\
        OpenCLIP ViT-B-32 (Zero-Shot) & 55.40 & -12.80 & 74.40 \\
        \bottomrule
    \end{tabular}
\end{table}

\begin{figure}[H]
    \centering
    \safeincludegraphics[width=0.85\linewidth]{figures/task1/patch_shuffle_examples.png}
    \caption{\textbf{Patch-Shuffled Visualizations and Model Predictions}: Examples showing surviving local evidence vs. destroyed global spatial arrangement.}
    \label{fig:task1_patch_examples}
\end{figure}
For patch shuffling, the results in Table~\ref{tab:task1_patch_shuffle} show drops in accuracy across all models, with OpenCLIP Probe showing the largest accuracy drop ($-16.80\%$) and Zero-Shot dropping to $55.40\%$ accuracy ($74.40\%$ consistency). In contrast, ResNet-50 and ViT-B/16 show only modest drops ($-6.00\%$ and $-7.20\%$) while retaining $>90\%$ accuracy and $>91\%$ consistency. This demonstrates that ResNet and ViT function effectively as bag-of-features classifiers: they retain high accuracy when global layout is destroyed because surviving local evidence inside individual $56 \times 56$ tiles is sufficient for classification (Figure~\ref{fig:task1_patch_examples}). OpenCLIP relies much more heavily on global scene composition learned from multimodal web pretraining, causing it to degrade more sharply when spatial arrangement is permuted.

\paragraph{Representation Geometry vs. Prediction Stability}
\begin{table}[H]
    \centering
    \caption{\textbf{Feature Representation Cosine Stability ($I_T$)}: Mean cosine similarity between clean representations and representations under each intervention.}
    \label{tab:task1_cosine_stability}
    \begin{tabular}{lcccc}
        \toprule
        \textbf{Model / Backbone} & \textbf{Grayscale} & \textbf{Cue Conflict} & \textbf{Translation ($\delta=32$)} & \textbf{Patch Shuffling} \\
        \midrule
        ResNet-50 & 0.8029 & 0.3651 & 0.9377 & 0.6953 \\
        ViT-B/16 & 0.7405 & 0.3701 & 0.9473 & 0.6333 \\
        OpenCLIP ViT-B-32 & 0.8891 & 0.7522 & 0.9597 & 0.7560 \\
        \bottomrule
    \end{tabular}
\end{table}

\begin{figure}[H]
    \centering
    \begin{subfigure}[b]{0.32\linewidth}
        \centering
        \safeincludegraphics[width=\linewidth]{figures/task1/umap_resnet50.png}
        \caption{ResNet-50 Projection}
        \label{fig:umap_resnet}
    \end{subfigure}
    \hfill
    \begin{subfigure}[b]{0.32\linewidth}
        \centering
        \safeincludegraphics[width=\linewidth]{figures/task1/umap_vitb16.png}
        \caption{ViT-B/16 Projection}
        \label{fig:umap_vit}
    \end{subfigure}
    \hfill
    \begin{subfigure}[b]{0.32\linewidth}
        \centering
        \safeincludegraphics[width=\linewidth]{figures/task1/umap_clip.png}
        \caption{OpenCLIP ViT-B-32 Projection}
        \label{fig:umap_clip}
    \end{subfigure}
    \caption{\textbf{2D Manifold Visualizations (Clean vs. Transformed)}: 2D UMAP projections of backbone representations combining clean and intervened examples. Points are colored by ground-truth class with circles denoting clean images and crosses denoting transformed images.}
    \label{fig:task1_manifold_projections}
\end{figure}
The results show that all interventions altered representations to varying degrees, revealing a clear mismatch between prediction stability and representation stability. For instance, while Grayscale caused only a minimal accuracy drop in Task 1.2 ($<2.5\%$), cosine stability dropped significantly to $0.7405$ in ViT and $0.8029$ in ResNet, showing that internal representations shifted substantially even though the final classifier output remained unchanged. Translation had the most stable representation ($I_T > 0.93$ across all models), whereas Cue Conflict and Patch Shuffling caused large disruptions, with Cue Conflict collapsing cosine stability to $\approx 0.37$ for ResNet and ViT. In the UMAP projections (Figure~\ref{fig:task1_manifold_projections}), clean and transformed points still cluster by class (explaining why accuracy remains high), but the transformed points are scattered and shifted within the clusters. Comparing OpenCLIP's zero-shot decisions with its trained linear head further confirms this: the trained linear head is able to navigate these representation shifts and achieve $97.40\%$ accuracy, whereas zero-shot cosine similarity with a fixed text prompt struggles under distribution drift.

%%%%%%%%%%%%%%%%%%%%%%%%%%%%%%%%%%%%%%%%%%%%%%%%%%%%%%%%%%%%%%%%%%%%%%%%%%%%%%%
\newpage
\section{Task 2: Unsupervised Domain Adaptation (UDA)}
\label{sec:task2}

\subsection{Methodology}
For our Unsupervised Domain Adaptation experiments, we made use of the PACS dataset, which consists of seven classes across four domains (Painting, Art, Cartoon, Sketch). We used Sketch as our unlabeled target domain. For each source domain, we used an 80/20 training/validation split with a seed of 6304.

We evaluated four distinct approaches: Source-Only (ERM), Mathematical Alignment using Deep Adaptation Network (DAN), Adversarial Alignment using Domain Adversarial Neural Network (DANN), and Class-Aware Adversarial Alignment using Conditional DANN (CDAN).

For our model, we used a fine-tuned ResNet-18 (\texttt{IMAGENET1K\_V1}) with a re-trained seven-class linear classifier head. The images were resized to $256 \times 256$ with a random $224 \times 224$ crop during training and a $224 \times 224$ center-crop for validation. We froze all BatchNorm running means and variances at their ImageNet values. We used the AdamW optimizer with a learning rate of $10^{-4}$ and weight decay of $10^{-4}$, stopping after five epochs without improvement in the mean source-validation macro-F1.

In order to perform DAN-style alignment, we added a 512-dimensional Maximum Mean Discrepancy (MMD) penalty to the feature space: $L_{DAN} = L_{\text{cls}} + \lambda_{\text{MMD}} \left\| E_s[\phi(F(x_s))] - E_t[\phi(F(x_t))] \right\|_2$, with $\lambda_{\text{MMD}}=1$ and a sum of three RBF kernels whose bandwidths were 0.5, 1, and 2 times the median pairwise squared feature distance in the current combined batch.

Next, a 256-unit discriminator with a dropout of 0.5, ReLU activation, and a two-class output layer was attached with the standard schedule $\alpha(p) = \frac{2}{1 + \exp(-10p)} - 1$.

For the CDAN implementation, the domain discriminator was given both the feature $f = F(x)$ and the classifier's probability vector $p = C(f)$ through $g(x) = \text{vec}(f \otimes p)$; the activation function, loss weight, dropout, and hidden width were kept exactly as in the DANN implementation.

The source-validation, target accuracy, macro-F1 scores, and domain separability scores (via a 70/30 logistic regression diagnostic) were then evaluated. In order to examine the relationship between domain invariance and target performance, we varied $\lambda_{\text{MMD}}$ across $\{0.1, 1.0, 10.0\}$ while keeping all other hyperparameters fixed.

\subsection{Results}

\paragraph{Training Dynamics and Convergence}
\begin{figure}[H]
    \centering
    \begin{subfigure}[b]{0.48\linewidth}
        \centering
        \safeincludegraphics[width=\linewidth]{figures/task2/training_loss_curves.png}
        \caption{Source Classification Loss}
        \label{fig:task2_cls_loss}
    \end{subfigure}
    \hfill
    \begin{subfigure}[b]{0.48\linewidth}
        \centering
        \safeincludegraphics[width=\linewidth]{figures/task2/alignment_loss_curves.png}
        \caption{Domain Alignment / Discriminator Loss}
        \label{fig:task2_align_loss}
    \end{subfigure}
    \caption{\textbf{Task 2 Training Dynamics}: Convergence curves showing (a) source classification cross-entropy and (b) alignment objectives (MMD penalty / domain discriminator loss) across adaptation epochs.}
\end{figure}
During training, the source classification loss for all four methods decreases steadily (Figure~\ref{fig:task2_cls_loss}), suggesting that the models learn the source features correctly without optimization difficulties. For the alignment objectives (Figure~\ref{fig:task2_align_loss}), DAN's MMD loss converges to $\sim 0.11$, while the adversarial (DANN) and conditional adversarial (CDAN) alignment methods stabilize at a higher loss around $0.69\text{--}0.73$. This is close to the theoretical minimax equilibrium of $-\ln(0.5) \approx 0.693$, suggesting that the adversarial training converged as designed.


\paragraph{Benchmark Comparison and Negative Transfer}
\begin{table}[H]
    \centering
    \caption{\textbf{Task 2 Benchmark Comparison on PACS (Target: Sketch)}: Source validation metrics per domain, mean source performance, target performance, accuracy gain relative to Source-only, and domain separability score (70/30 logistic regression; 50\% = chance).}
    \label{tab:task2_main_results}
    \resizebox{\linewidth}{!}{%
    \setlength{\tabcolsep}{4pt}
    \begin{tabular}{lccccccccc}
        \toprule
        & \multicolumn{3}{c}{\textbf{Source Val Accuracy (\%)}} & \textbf{Mean Source} & \textbf{Mean Source} & \textbf{Target} & \textbf{Target} & \textbf{Target Gain} & \textbf{Domain Sep.} \\
        \cmidrule(lr){2-4}
        \textbf{Method} & \textbf{Photo} & \textbf{Art} & \textbf{Cartoon} & \textbf{Accuracy (\%)} & \textbf{Macro-F1} & \textbf{Acc (\%)} & \textbf{F1} & \textbf{$\Delta$ Acc (\%)} & \textbf{Score (\%)} \\
        \midrule
        Source-only & 95.81 & 90.49 & 95.10 & 93.80 & 0.94 & 72.74 & 0.75 & +0.00 & 99.59 \\
        DAN & 96.71 & 91.46 & 94.88 & 94.35 & 0.94 & 63.86 & 0.56 & -8.88 & 90.93 \\
        DANN & 96.11 & 91.71 & 94.46 & 94.09 & 0.94 & 19.65 & 0.05 & -53.09 & 100.00 \\
        CDAN & 95.51 & 90.00 & 94.88 & 93.46 & 0.93 & 2.04 & 0.01 & -70.71 & 100.00 \\
        \bottomrule
    \end{tabular}%
    }
\end{table}

\begin{table}[H]
    \centering
    \caption{\textbf{Per-Class Target Accuracy (Sketch)}: Class-level breakdown across adaptation methods to inspect negative transfer.}
    \label{tab:task2_per_class}
    \begin{tabular}{lccccccc}
        \toprule
        \textbf{Method} & \textbf{Dog} & \textbf{Elephant} & \textbf{Giraffe} & \textbf{Guitar} & \textbf{Horse} & \textbf{House} & \textbf{Person} \\
        \midrule
        Source-only & 57.77 & 83.38 & 57.90 & 96.88 & 71.08 & 98.75 & 69.38 \\
        DAN & 50.91 & 72.30 & 61.35 & 84.38 & 64.46 & 100.00 & 0.00 \\
        DANN & 100.00 & 0.00 & 0.00 & 0.00 & 0.00 & 0.00 & 0.00 \\
        CDAN & 0.00 & 0.00 & 0.00 & 0.00 & 0.00 & 100.00 & 0.00 \\
        \bottomrule
    \end{tabular}
\end{table}

\begin{figure}[H]
    \centering
    \safeincludegraphics[width=0.9\linewidth]{figures/task2/confusion_failures.png}
    \caption{\textbf{Target Error Analysis on Sketch}: Dominant confusions and failure cases illustrating semantic preservation vs. negative transfer across methods.}
    \label{fig:task2_confusions}
\end{figure}
While all four approaches yield high mean accuracies on the source domains ($93.46\%\text{--}94.35\%$), all three adaptation methods suffered negative transfer relative to Source-only ($72.74\%$), with target performance differing greatly depending on the alignment strategy (Table~\ref{tab:task2_main_results}). Adversarial alignments (DANN and CDAN) resulted in catastrophic drops in target accuracy ($19.65\%$ and $2.04\%$) alongside near-perfect domain separability scores ($100.00\%$). Even for DAN, where domain separability was successfully reduced to $90.93\%$, target accuracy degraded by $-8.88\%$; therefore, a link between lower domain separability and higher target accuracy could not be established. Table~\ref{tab:task2_per_class} and Figure~\ref{fig:task2_confusions} provide further insight into these failures: adversarial methods suffered mode collapse, with DANN predicting `dog' $100\%$ of the time and CDAN predicting `house' $100\%$ of the time. While the per-class target accuracies for DAN were much more uniform across most categories, it completely failed to predict the `person' class ($0.00\%$). It is therefore important to analyse these aggregate metrics within the context of the class-specific results since they can conceal vital weaknesses of the approach, such as DAN's failure to predict the `person' class.


\paragraph{Controlled Alignment Strength Study}
\begin{table}[H]
    \centering
    \caption{\textbf{Controlled Alignment Strength Study}: Impact of alignment pressure on source performance, domain separability, and target recognition.}
    \label{tab:task2_design_study}
    \begin{tabular}{lccccc}
        \toprule
        \textbf{Hyperparameter Value} & \textbf{Mean Source Acc (\%)} & \textbf{Mean Source F1} & \textbf{Domain Sep. (\%)} & \textbf{Target Acc (\%)} & \textbf{Target F1} \\
        \midrule
        Setting 1 ($\lambda=0.1$) & 94.80 & 0.95 & 98.08 & 71.90 & 0.70 \\
        Setting 2 ($\lambda=1.0$) & 93.10 & 0.93 & 87.91 & 69.33 & 0.63 \\
        Setting 3 ($\lambda=10.0$) & 90.51 & 0.90 & 85.71 & 64.29 & 0.58 \\
        \bottomrule
    \end{tabular}
\end{table}
Increasing alignment pressure $\lambda_{\text{MMD}}$ from $0.1 \rightarrow 10.0$ decreased domain separability along with target and source accuracy (Table~\ref{tab:task2_design_study}), with $\lambda=0.1$ yielding the highest source accuracy ($94.80\%$) and F1 score ($0.95$). This indicates that while increasing alignment pressure decreases domain separability (increasing domain invariance), it comes at the expense of decreased class separability.


%%%%%%%%%%%%%%%%%%%%%%%%%%%%%%%%%%%%%%%%%%%%%%%%%%%%%%%%%%%%%%%%%%%%%%%%%%%%%%%
\newpage
\section{Task 3: Domain Generalization (DG)}
\label{sec:task3}

\subsection{Methodology}
Under the Domain Generalization setup, the target domain (Sketch) was strictly unseen during all stages of training, validation, model selection, and hyperparameter tuning. We reused the shared PACS splits with an 80/20 train/validation split (seed 6304), fine-tuning a ResNet-18 model initialized with ImageNet weights and keeping BatchNorm running statistics frozen.

For Domain Generalization, we evaluated two distinct paradigms: Pairwise Source-Domain Alignment and Parameter-Space Stability.

We reused the Source-only ERM baseline from the previous experiment trained on the three source domains (Photo, Art Painting, Cartoon). Its objective gives every source domain equal influence through domain-balanced batches:
\begin{equation}
L_{\text{ERM}}(\theta) = \frac{1}{3} \sum_{e \in \{P, A, C\}} R_e(\theta),
\end{equation}
where $P$, $A$, and $C$ denote Photo, Art Painting, and Cartoon. We calculated validation accuracy and macro-F1 scores across source domains.

Next, for Pairwise Source-Domain Alignment (DAN-DG), a new model was trained using the ERM classification loss plus the average MMD discrepancy over the three unordered pairs of source domains:
\begin{equation}
L_{\text{DAN-DG}}(\theta) = L_{\text{ERM}}(\theta) + \lambda_{\text{DG}} \sum_{(e, e')} \text{MMD}^2(F(X_e), F(X_{e'})).
\end{equation}
We applied MMD to the 512-dimensional feature representations immediately before the classifier head with $\lambda_{\text{DG}} = 1.0$ and a sum of three RBF kernels whose bandwidths were 0.5, 1.0, and 2.0 times the median pairwise squared feature distance in the current combined batch.

For Parameter-Space Stability, we trained the ERM objective with standard, non-adaptive Sharpness-Aware Minimization (SAM):
\begin{equation}
\min_\theta \max_{\|\epsilon\|_2 \le \rho} L_{\text{ERM}}(\theta + \epsilon),
\end{equation}
with perturbation radius $\rho = 0.05$, using the AdamW optimizer with learning rate $10^{-4}$, weight decay $10^{-4}$, and the same frozen BatchNorm running-statistics policy.

Source validation accuracy, macro-F1, worst-source validation accuracy, and final accuracy and macro-F1 on unseen Sketch were then reported.

To measure how much domain information remained in each representation, we froze the backbone, collected balanced features from the three source validation sets, created a 70/30 split with seed 6304, and trained a multinomial logistic-regression classifier with $C = 1.0$ to predict Photo, Art Painting, or Cartoon. 

We also used the same standardized local sharpness proxy for all three models, measuring the increase in validation cross-entropy loss after one normalized gradient-ascent perturbation with radius $\rho = 0.05$:
\begin{equation}
\Delta_{\text{sharp}} = L_{\text{val}}(\theta + \epsilon) - L_{\text{val}}(\theta), \quad \text{where } \epsilon = 0.05 \frac{\nabla_\theta L_{\text{val}}}{\|\nabla_\theta L_{\text{val}}\|_2}.
\end{equation}

\subsection{Results}

\paragraph{Domain Generalization Benchmark Performance}
\begin{table}[H]
    \centering
    \caption{\textbf{Task 3 Domain Generalization Performance on PACS}: Comparison of ERM, DAN-DG, and SAM across source validation domains, worst-source domain, unseen target (Sketch), and diagnostic metrics.}
    \label{tab:task3_main_results}
    \resizebox{\linewidth}{!}{%
    \setlength{\tabcolsep}{4pt}
    \begin{tabular}{lccccccccc}
        \toprule
        & \multicolumn{3}{c}{\textbf{Source Val Acc (\%)}} & \textbf{Mean Source} & \textbf{Worst Source} & \textbf{Sketch Target} & \textbf{Sketch} & \textbf{Source Domain} & \textbf{Sharpness Proxy} \\
        \cmidrule(lr){2-4}
        \textbf{Method} & \textbf{Photo} & \textbf{Art} & \textbf{Cartoon} & \textbf{Acc (\%)} & \textbf{Acc (\%)} & \textbf{Acc (\%)} & \textbf{$\Delta$ Acc (\%)} & \textbf{Sep. Score (\%)} & \textbf{$\Delta_{\text{sharp}}$} \\
        \midrule
        ERM & 95.81 & 90.49 & 95.10 & 93.80 & 90.49 & 72.74 & +0.00 & 87.04 & 0.71 \\
        DAN-DG & 96.11 & 92.20 & 94.03 & 94.11 & 92.20 & 61.85 & -10.89 & 73.09 & 0.27 \\
        SAM & 97.60 & 92.44 & 95.74 & 95.26 & 92.44 & 72.00 & -0.74 & 88.70 & 0.12 \\
        \bottomrule
    \end{tabular}%
    }
\end{table}

\begin{figure}[H]
    \centering
    \begin{subfigure}[b]{0.48\linewidth}
        \centering
        \safeincludegraphics[width=\linewidth]{figures/task3/dg_loss_curves.png}
        \caption{Classification \& MMD Penalty Curves}
        \label{fig:task3_loss}
    \end{subfigure}
    \hfill
    \begin{subfigure}[b]{0.48\linewidth}
        \centering
        \safeincludegraphics[width=\linewidth]{figures/task3/sharpness_vs_generalization.png}
        \caption{Sharpness Proxy vs. Sketch Accuracy}
        \label{fig:task3_sharpness_scatter}
    \end{subfigure}
    \caption{\textbf{Domain Generalization Dynamics and Diagnostics}: (a) Training loss and pairwise MMD penalty trajectories; (b) Empirical correlation between the standardized local sharpness proxy $\Delta_{\text{sharp}}$ and unseen Sketch target performance.}
    \label{fig:task3_diagnostics}
\end{figure}
Both DAN-DG and SAM achieved higher mean source validation accuracies ($94.11\%$ and $95.26\%$) and higher worst-source accuracies ($92.20\%$ and $92.44\%$) than ERM ($93.80\%$ and $90.49\%$), but both resulted in lower unseen Sketch target accuracy (Table~\ref{tab:task3_main_results}). SAM resulted in the highest mean and worst-source performance while retaining competitive target accuracy ($72.00\%$ vs. $72.74\%$ for ERM). These results demonstrate that higher source validation performance does not reliably translate into improved unseen domain generalization.

Furthermore, DAN-DG resulted in substantially lower source domain separability ($73.09\%$) compared to ERM ($87.04\%$) and SAM ($88.70\%$), confirming that pairwise MMD successfully made observed source domains less distinguishable. However, this cross-source invariance severely impaired target recognition ($61.85\%$), indicating that forcing domain alignment stripped away class-discriminative cues needed to classify sketches.

Regarding parameter stability, Figure~\ref{fig:task3_diagnostics}(a) shows steady convergence in classification loss and pairwise MMD penalty. SAM successfully converged to a much flatter loss minimum, achieving the lowest standardized sharpness proxy ($\Delta_{\text{sharp}} = 0.12$ compared to $0.71$ for ERM). Crucially, as visualized in Figure~\ref{fig:task3_sharpness_scatter}, our experiments did not find a clear relationship between local stability and sketch performance. While SAM had the highest stability ranking by a large margin, it had the second highest sketch accuracy and was only slightly lower than the method with the lowest stability (ERM). Moreover, while DAN-DG had significantly higher local stability than ERM, it resulted in the lowest target accuracy ($61.85\%$).


\paragraph{Class-Level Transfer and Controlled Perturbation Study}
\begin{table}[H]
    \centering
    \caption{\textbf{Per-Class Accuracy on Unseen Sketch}: Comparison between DG methods (Task 3) and UDA methods with target access (Task 2).}
    \label{tab:task3_per_class}
    \begin{tabular}{lccccccc}
        \toprule
        \textbf{Method} & \textbf{Dog} & \textbf{Elephant} & \textbf{Giraffe} & \textbf{Guitar} & \textbf{Horse} & \textbf{House} & \textbf{Person} \\
        \midrule
        ERM (Shared) & 57.77 & 83.38 & 57.90 & 96.88 & 71.08 & 98.75 & 69.38 \\
        DAN-DG (No Target) & 63.60 & 81.76 & 30.68 & 82.89 & 47.06 & 87.50 & 90.62 \\
        DAN (UDA with Target) & 50.91 & 72.30 & 61.35 & 84.38 & 64.46 & 100.00 & 0.00 \\
        SAM (No Target) & 28.50 & 94.59 & 63.21 & 87.17 & 87.01 & 76.25 & 82.50 \\
        \bottomrule
    \end{tabular}
\end{table}

\begin{table}[H]
    \centering
    \caption{\textbf{Controlled DG Study (SAM)}: Effect of varying perturbation radius ($\rho$) on source performance, diagnostics, and unseen Sketch recognition.}
    \label{tab:task3_design_study}
    \begin{tabular}{lccccc}
        \toprule
        \textbf{Hyperparameter Value} & \textbf{Mean Source Acc (\%)} & \textbf{Worst Source Acc (\%)} & \textbf{Diagnostic Metric} & \textbf{Sketch Acc (\%)} & \textbf{Sketch F1} \\
        \midrule
        $\rho = 0.05$ (Nominal) & 95.26 & 92.44 & 0.12 & 72.00 & 0.72 \\
        $\rho = 0.01$ & 94.76 & 91.71 & 0.17 & 75.74 & 0.77 \\
        $\rho = 0.1$ & 93.69 & 89.02 & 0.10 & 62.99 & 0.61 \\
        \bottomrule
    \end{tabular}
\end{table}
Analyzing the class-based accuracies of the approaches reveals noteworthy details which could not be expressed by aggregate metrics alone (Table~\ref{tab:task3_per_class}). While DAN-DG has a lower overall accuracy compared to DAN, its accuracies are more uniform. DAN-DG attained a 90.62\% accuracy compared to 0\% for DAN. Though it is important to note that while DAN-DG had no absolute class failures it did perform weakly on some classes such as Giraffe and Horse. Access to unlabeled sketches did not result in superior performance to target-free methods. Both SAM and ERM outperformed target-aware DAN significantly and access to unlabelled sketches resulted in only a 2\% increase in accuracy.

A lower perturbation radius of $\rho = 0.01$ resulted in higher sketch accuracies and F1-scores, while a radius of $0.05$ gave the best source accuracy (Table~\ref{tab:task3_design_study}). Interestingly, the source accuracy decreased when the radius was increased further to $0.1$. The results indicate an inverse relationship between perturbation radius and sharpness proxy as well as perturbation radius and sketch accuracy/F1 score, but no clear effect on source accuracy.


%%%%%%%%%%%%%%%%%%%%%%%%%%%%%%%%%%%%%%%%%%%%%%%%%%%%%%%%%%%%%%%%%%%%%%%%%%%%%%%
\newpage
\section{Task 4: Open-Set Recognition (OSR)}
\label{sec:task4}

\subsection{Methodology}
To evaluate how classifiers handle out-of-distribution inputs, we used CIFAR-10 (10 classes with a stratified 90/10 train/validation split using seed 6304 and the official 10,000-image test set) to represent the known classes. To construct semantic unknowns, we extracted two disjoint 800-image subsets from the CIFAR-100 test set:
\begin{enumerate}
    \item \textbf{Near Unknowns} (semantically similar categories): \texttt{bus}, \texttt{pickup\_truck}, \texttt{motorcycle}, \texttt{tractor}, \texttt{wolf}, \texttt{fox}, \texttt{leopard}, and \texttt{camel}.
    \item \textbf{Far Unknowns} (semantically distant categories): \texttt{bottle}, \texttt{bowl}, \texttt{chair}, \texttt{clock}, \texttt{keyboard}, \texttt{mushroom}, \texttt{sunflower}, and \texttt{wardrobe}.
\end{enumerate}

For our backbone, we used a CIFAR-adapted ResNet-18 where the initial $7 \times 7$ convolution was replaced with a $3 \times 3$ kernel (stride 1) and the max-pooling layer was removed to preserve $32 \times 32$ spatial resolution. The vanilla model was trained using SGD with momentum 0.9, an initial learning rate of 0.1 decayed via a cosine schedule, weight decay of $5 \times 10^{-4}$, batch size 128, and seed 6304 for 100 epochs.

We evaluated four post-hoc novelty scoring functions $u(x)$ where higher scores indicate greater unknownness:
\begin{enumerate}
    \item \textbf{Maximum Softmax Probability (MSP):} $u_{\text{MSP}}(x) = 1 - \max_k p_k(x)$, where $p = \text{Softmax}(z(x))$.
    \item \textbf{Maximum Logit Score (MLS):} $u_{\text{MLS}}(x) = -\max_k z_k(x)$.
    \item \textbf{Energy Score:} $u_{\text{Energy}}(x) = -\log \sum_k \exp(z_k(x))$.
    \item \textbf{Mahalanobis Distance:} $u_{\text{Mah}}(x) = \min_c (f(x) - \mu_c)^T \Sigma^{-1} (f(x) - \mu_c)$, using class centroids $\mu_c$ and a shared diagonal covariance $\Sigma$ computed over unaugmented training features.
\end{enumerate}
Decision thresholds $\tau$ were calibrated at the $95^{\text{th}}$ percentile of known validation unknownness scores to enforce a fixed 95\% True Positive Rate (5\% false rejection rate) on knowns.

Beyond post-hoc scoring, we evaluated two trained-model paradigms:
\begin{itemize}
    \item \textbf{GCSC (Good Closed-Set Classifier):} Trains the modified ResNet-18 with aggressive RandAugment (2 operations, magnitude 9) under the vanilla recipe to test if stronger closed-set representation alone improves open-set detection.
    \item \textbf{PROSER (Learning Placeholders):} Adds dummy classifier prototypes and dummy data placeholders generated via feature mixup between second-most-likely classes to explicitly carve out an unknown decision boundary. The backbone was fine-tuned for 50 epochs using SGD with learning rate $10^{-3}$, momentum 0.9, weight decay $5 \times 10^{-4}$, and cosine decay.
\end{itemize}

Models were evaluated on Closed-Set Accuracy (CSA), AUROC across Near, Far, and All Unknowns, and false positive rejection rates at $\tau$. Finally, for qualitative failure analysis, unknown images accepted under the vanilla MLS threshold ($u_{\text{MLS}}(x) \le \tau$) were extracted to inspect semantic overlap and feature aliasing.

\subsection{Results}

\paragraph{Post-Hoc Novelty Scoring Evaluation}
\begin{table}[H]
    \centering
    \caption{\textbf{Post-Hoc Novelty Detection Metrics on Vanilla ResNet-18}: AUROC across unknown subsets and rejection performance calibrated at $\tau = 95^{\text{th}}$ percentile of CIFAR-10 validation unknownness.}
    \label{tab:task4_posthoc_scores}
    \begin{tabular}{lcccccc}
        \toprule
        & \multicolumn{3}{c}{\textbf{AUROC (\%)}} & \multicolumn{3}{c}{\textbf{Rejection Performance at 95\% Val TPR}} \\
        \cmidrule(lr){2-4} \cmidrule(lr){5-7}
        \textbf{Score} & \textbf{Near} & \textbf{Far} & \textbf{All} & \textbf{CIFAR-10 Test Acc. (\%)} & \textbf{Near Rej. (\%)} & \textbf{Far Rej. (\%)} \\
        \midrule
        MSP & 80.84 & 89.32 & 85.08 & 94.80 & 29.38 & 47.62 \\
        MLS & 79.36 & 88.47 & 83.91 & 94.80 & 32.88 & 54.25 \\
        Energy & 79.35 & 88.49 & 83.92 & 94.80 & 33.38 & 54.00 \\
        Mahalanobis & 77.74 & 91.33 & 84.53 & 94.80 & 30.00 & 53.88 \\
        \bottomrule
    \end{tabular}
\end{table}

\begin{figure}[H]
    \centering
    \begin{subfigure}[b]{0.48\linewidth}
        \centering
        \safeincludegraphics[width=\linewidth]{figures/task4/score_distributions.png}
        \caption{Unknownness Score Distributions}
        \label{fig:task4_score_dist}
    \end{subfigure}
    \hfill
    \begin{subfigure}[b]{0.48\linewidth}
        \centering
        \safeincludegraphics[width=\linewidth]{figures/task4/roc_curves.png}
        \caption{ROC Curves (Known vs. All Unknowns)}
        \label{fig:task4_roc_curves}
    \end{subfigure}
    \caption{\textbf{Novelty Score Separation}: (a) Density distributions for knowns vs. near and far unknowns; (b) ROC curves across novelty scoring functions.}
    \label{fig:task4_scoring_analysis}
\end{figure}
Across all four metrics, far unknowns were significantly easier to detect and reject than near unknowns, with AUROC being approximately $10\%$ higher for far unknowns compared to near unknowns, and Rejection (TPR) being $15\text{--}25\%$ higher (Table~\ref{tab:task4_posthoc_scores}). MSP achieved the highest overall AUROC, while Mahalanobis performed the best out of the four for far unknowns. MLS and Energy provided superior thresholds for rejecting both far and near unknowns. The results indicate that feature distance can be most informative when unknowns are far from known clusters since Mahalanobis yields higher scores for far unknowns, while unnormalized logit scores (MLS and Energy) are more informative for establishing rejection thresholds.

Visualizing the score distributions (Figure~\ref{fig:task4_scoring_analysis}) shows that there is significant overlap between all three distributions, but the novelty score is consistently increasing for both far and near unknowns.

\paragraph{Trained-Model OSR Benchmark and Failure Analysis}
\begin{table}[H]
    \centering
    \caption{\textbf{Trained-Model OSR Benchmark}: Comparison of Vanilla, GCSC, and PROSER on Closed-Set Accuracy (CSA), AUROC, and validation-calibrated rejection rates at 95\% TPR.}
    \label{tab:task4_trained_models}
    \footnotesize
    \begin{tabular}{lccccccc}
        \toprule
        & & \multicolumn{3}{c}{\textbf{AUROC (\%)}} & \multicolumn{3}{c}{\textbf{Rejection Performance at 95\% Val TPR}} \\
        \cmidrule(lr){3-5} \cmidrule(lr){6-8}
        \textbf{Model / Score} & \textbf{CSA (\%)} & \textbf{Near} & \textbf{Far} & \textbf{All} & \textbf{FPR@95TPR (\%)} & \textbf{Near Rej. (\%)} & \textbf{Far Rej. (\%)} \\
        \midrule
        Vanilla (MLS) & 94.80 & 79.36 & 88.47 & 83.91 & 5.40 & 32.88 & 54.25 \\
        GCSC (MLS) & 95.37 & 81.82 & 90.40 & 86.11 & 5.35 & 37.00 & 61.12 \\
        PROSER (MLS) & 94.71 & 79.92 & 87.42 & 83.67 & 5.06 & 30.75 & 49.38 \\
        PROSER (Score) & 94.71 & 74.71 & 89.91 & 82.31 & 5.18 & 28.62 & 56.38 \\
        \bottomrule
    \end{tabular}
\end{table}

\begin{figure}[H]
    \centering
    \begin{subfigure}[b]{0.48\linewidth}
        \centering
        \safeincludegraphics[width=\linewidth]{figures/task4/near_unknown_failures.png}
        \caption{Accepted Near Unknowns}
        \label{fig:near_failures}
    \end{subfigure}
    \hfill
    \begin{subfigure}[b]{0.48\linewidth}
        \centering
        \safeincludegraphics[width=\linewidth]{figures/task4/far_unknown_failures.png}
        \caption{Accepted Far Unknowns}
        \label{fig:far_failures}
    \end{subfigure}
    \caption{\textbf{Failure Inspection Under Vanilla MLS Threshold $\tau$}: Incorrectly accepted unknown inputs showing unknown class, predicted CIFAR-10 class, logit score, and margin relative to threshold.}
    \label{fig:task4_failures}
\end{figure}

\begin{table}[H]
    \centering
    \caption{\textbf{Case Studies of Accepted Unknown Inputs}: Inspection of three near-unknown and three far-unknown failures under the vanilla MLS decision boundary.}
    \label{tab:task4_failure_cases}
    \begin{tabular}{lcccc}
        \toprule
        \textbf{Type} & \textbf{True CIFAR-100 Class} & \textbf{Predicted CIFAR-10 Class} & \textbf{$u_{\text{MLS}}(x)$} & \textbf{Threshold $\tau$} \\
        \midrule
        Near & pickup\_truck & automobile & -11.22 & -5.85 \\
        Near & bus & truck & -12.05 & -5.85 \\
        Near & wolf & bird & -11.36 & -5.85 \\
        \midrule
        Far & wardrobe & truck & -10.60 & -5.85 \\
        Far & mushroom & bird & -11.03 & -5.85 \\
        Far & chair & airplane & -10.08 & -5.85 \\
        \bottomrule
    \end{tabular}
\end{table}
As shown in Table~\ref{tab:task4_trained_models}, GCSC attained the highest closed-set accuracy, highest AUROC across all splits, and highest rejection rates, followed by Vanilla ResNet-18. The closed-set accuracy and AUROC were both higher for GCSC; this could imply that RandAugment tends to improve closed-set accuracy and open-set recognition. PROSER had a closed-set accuracy of 94.71\%. When comparing MLS to its learned unknown placeholder score, Near AUROC and Near rejection fell while Far AUROC and Far rejection rose, indicating that using an unknown placeholder score can help with detecting and rejecting far unknowns. Interestingly, PROSER did not yield higher rejection performance scores consistently, though it did have better performance on far unknowns compared to vanilla.

Looking at the failure cases under the vanilla MLS threshold (Figure~\ref{fig:task4_failures} and Table~\ref{tab:task4_failure_cases}), the most frequently accepted near unknowns were pickup truck (94\%), bus (84\%), and wolf (79\%), which were absorbed by truck, automobile, and cat/dog. These failures are semantically plausible due to shared visual and category overlap. We do not see the same trend for far unknowns however, where the most accepted were wardrobe (69\%), mushroom (48\%), and chair (47\%) absorbed by truck/cat, bird, and airplane, so semantic overlap alone is insufficient to explain the failures. Feature aliasing could be a plausible explanation for the far unknown failures though the results in this paper are insufficient to establish an exact causal relationship for these failures.

%%%%%%%%%%%%%%%%%%%%%%%%%%%%%%%%%%%%%%%%%%%%%%%%%%%%%%%%%%%%%%%%%%%%%%%%%%%%%%%
\newpage
\section{Discussion (cross-task)}
\label{sec:discussion}

\subsection{Architectural Bias vs. Pretraining, Supervision, and Capacity}
Since OpenCLIP and ViT both have Transformer-based architectures, some of their differences can adequately be explained by differences in pretraining. It is reasonable to conclude from our results that OpenCLIP's high cosine stability across interventions ($I_T \ge 0.75$) is likely because it was pretrained with InfoNCE loss over diverse web images, which trains the model to keep matching image and text representations aligned and stable. ViT and OpenCLIP's higher shape bias likely stems from architectural differences between CNNs and Transformers: the evidence for this is that the pretraining for ViT and ResNet is similar (both being trained on ImageNet with cross-entropy loss), yet ViT achieved a much higher shape bias ($94.22\%$ vs. $78.81\%$). Furthermore, ViT's higher coverage ($78.64\%$ vs. $68.64\%$) means that it would be inappropriate to attribute this difference to noise. In terms of supervision and classifier heads, OpenCLIP's low zero-shot performance ($68.20\%$) and high third-class confusion ($N_{\text{other}}=85$) largely disappeared once the linear probe was trained ($97.40\%$ accuracy, $90.80\%$ shape bias), proving this was an artifact of the static prompt bottleneck rather than the underlying representation. Finally, OpenCLIP's greater sensitivity to 32-pixel translation compared to ViT-B/16 reflects the architectural difference in patch token size ($32 \times 32$ vs. $16 \times 16$).

\subsection{Cross-Task Representation Geometry and Invariance}
Our experiments could not find a link between domain invariance and improved recognition. Increasing alignment pressure in our domain adaptation experiments yielded a lower target accuracy, and adversarial alignment methods suffered mode collapse. In our Domain Generalization experiments, Pairwise MMD performed the best at reducing domain separability but yielded the lowest target accuracy on Sketch images. 

Instead of improving recognition, our alignment interventions actually decreased accuracy. Though it is important to mention that these experiments alone cannot provide a causal explanation for this behaviour, one possible explanation could be that visual cues that are important for category recognition are removed or distorted during domain alignment, which hurts recognition on unseen domains.


\subsection{Target Access vs. Robust Optimization}
Comparing the accuracy of target-aware DAN from our UDA experiments ($63.86\%$) to our target-free DAN-DG from our Domain Generalization experiments ($61.85\%$) indicates that having access to unlabeled Sketch images did improve accuracy, but only slightly ($\sim 2\%$).

It is unclear whether the increase was caused by access to the target domain or by the difference in how the loss for the two approaches is computed. DAN uses source-to-target alignment and DAN-DG uses pairwise cross-source alignment, so it is plausible that the difference in accuracy could be a result of either of these factors. Moreover, DANN and CDAN both had access to the target data but suffered mode collapse, while SAM, which did not have access to target data, outperformed all target-aligned models. 

Therefore, while the data did show improved performance for some methods when they had access to sketch data, it is insufficient by itself to establish a causal relationship.


\subsection{Closed-Set Discrimination vs. Open-Set Rejection}
Vanilla ResNet had high test accuracy on CIFAR-10 but falsely accepted a large percentage of unknowns. Therefore, high closed-set accuracy does not guarantee strong open-set performance, because standard classifiers partition the entire feature space without an option to abstain.

GCSC's RandAugment improved closed-set accuracy ($95.37\%$) and rejection across all splits. On the other hand, reserving probability mass for placeholder classes in PROSER helped the model reject far unknowns, but at the cost of weaker performance on near unknowns and a slight drop in closed-set accuracy ($94.71\%$).

Across all approaches, far unknowns were significantly easier to both detect and reject compared to near unknowns, which is expected due to the semantic overlap between near unknowns and known data.

%%%%%%%%%%%%%%%%%%%%%%%%%%%%%%%%%%%%%%%%%%%%%%%%%%%%%%%%%%%%%%%%%%%%%%%%%%%%%%%
\section{Conclusion}
\label{sec:conclusion}

Based on evidence across all four experiments, we can conclude that a robust representation should preserve class-separating evidence (for example, local textures for ResNet and global shape for ViT). It should be invariant to superficial variations such as color and small translations which do not alter the label, but should not force feature distributions to overlap when domain style and class cues are entangled. Since finding flat minima in parameter space proved far more effective for generalization than increasing alignment pressure between domains, it is also ideal for the representation to be invariant to curvature in the loss landscape. Finally, inputs which fall outside the training support such as far unknowns should not be forced into a known class, and ambiguous near unknowns require explicit rejection thresholds to safely abstain.

\newpage
\bibliographystyle{plain}
\begin{thebibliography}{99}

\bibitem{geirhos2019imagenet}
R.~Geirhos, P.~Rubisch, C.~Michaelis, M.~Bethge, F.~A. Wichmann, and W.~Brendel.
\newblock Image{N}et-trained {CNN}s are biased towards texture; increasing shape bias improves accuracy and robustness.
\newblock In \emph{International Conference on Learning Representations (ICLR)}, 2019.

\bibitem{raghu2021vision}
M.~Raghu, T.~Unterthiner, S.~Kornblith, M.~Zhang, and A.~Dosovitskiy.
\newblock Do vision transformers see like convolutional neural networks?
\newblock In \emph{Advances in Neural Information Processing Systems (NeurIPS)}, volume~34, pages 12111--12121, 2021.

\bibitem{ganin2016domain}
Y.~Ganin, E.~Ustinova, H.~Ajakan, P.~Germain, H.~Larochelle, F.~Laviolette, M.~Marchand, and V.~Lempitsky.
\newblock Domain-adversarial training of neural networks.
\newblock \emph{Journal of Machine Learning Research}, 17(59):1--35, 2016.

\bibitem{long2015learning}
M.~Long, Y.~Cao, J.~Wang, and M.~Jordan.
\newblock Learning transferable features with deep adaptation networks.
\newblock In \emph{International Conference on Machine Learning (ICML)}, pages 97--105, 2015.

\bibitem{long2018conditional}
M.~Long, Z.~Cao, J.~Wang, and M.~I. Jordan.
\newblock Conditional adversarial domain adaptation.
\newblock In \emph{Advances in Neural Information Processing Systems (NeurIPS)}, 2018.

\bibitem{foret2021sharpness}
P.~Foret, A.~Kleiner, H.~Mobahi, and B.~Neyshabur.
\newblock Sharpness-aware minimization for efficiently improving generalization.
\newblock In \emph{International Conference on Learning Representations (ICLR)}, 2021.

\bibitem{vaze2022open}
S.~Vaze, K.~Han, A.~Vedaldi, and A.~Zisserman.
\newblock Open-set recognition: a good closed-set classifier is all you need?
\newblock In \emph{International Conference on Learning Representations (ICLR)}, 2022.

\bibitem{zhou2021learning}
D.-W. Zhou, H.-J. Ye, and D.-C. Zhan.
\newblock Learning placeholders for open-set recognition.
\newblock In \emph{IEEE/CVF Conference on Computer Vision and Pattern Recognition (CVPR)}, pages 4401--4410, 2021.

\bibitem{hendrycks2017baseline}
D.~Hendrycks and K.~Gimpel.
\newblock A baseline for detecting misclassified and out-of-distribution examples in neural networks.
\newblock In \emph{International Conference on Learning Representations (ICLR)}, 2017.

\bibitem{liu2020energy}
W.~Liu, X.~Wang, J.~Owens, and Y.~Li.
\newblock Energy-based out-of-distribution detection.
\newblock In \emph{Advances in Neural Information Processing Systems (NeurIPS)}, 2020.

\bibitem{lee2018simple}
K.~Lee, K.~Lee, H.~Lee, and J.~Shin.
\newblock A simple unified framework for detecting out-of-distribution samples and adversarial attacks.
\newblock In \emph{Advances in Neural Information Processing Systems (NeurIPS)}, 2018.

\bibitem{radford2021learning}
A.~Radford, J.~W. Kim, C.~Hallacy, A.~Ramesh, G.~Goh, S.~Agarwal, G.~Sastry, A.~Askell, P.~Mishkin, J.~Clark, et~al.
\newblock Learning transferable visual models from natural language supervision.
\newblock In \emph{International Conference on Machine Learning (ICML)}, pages 8748--8763, 2021.

\end{thebibliography}

\end{document}
